% ================================================================
% ch04_bhakti_hetu.tex — अध्याय ४: भक्ति हेतु
% ================================================================

\cleardoublepage
\section*{अध्याय ४: भक्ति हेतु}
\addcontentsline{toc}{section}{अध्याय ४: भक्ति हेतु}

\subsection*{४.१ भक्ति के लक्षण आ शूरवीरता}
\addcontentsline{toc}{subsection}{४.१ भक्ति के लक्षण आ शूरवीरता}

\textbf{मूल भोजपुरी रचना:}

\begin{quote}
\addfontfeatures{Color=3D3428}
भक्ति करे सो सुरमा, तन मन लज्जा खोय।\\
छैल चिकनिया बिसनी, वा से भक्ति न होय॥
\end{quote}

\textbf{भोजपुरी भावार्थ:} सच्चा भक्त उहे वीर ह जे तन, मन आ संसार के लोक-लाज छोड़ के प्रभु के शरण में आ जाला। छैल-चिकनिया आ भोगी-विलास में डूबल लोग से भक्ति ना हो सके।

\begin{sloppypar}
{\engfont \textit{English Translation:} True devotion requires real courage --- surrendering bodily pride and worldly shame.
The superficial and indulgence-seeking cannot tread the earnest path of devotion.}
\end{sloppypar}

\begin{center}
{\color{bordercolor}\rule{0.4\textwidth}{0.4pt}}
\end{center}

\subsection*{४.२ निष्काम सेवा आ अमर पद}
\addcontentsline{toc}{subsection}{४.२ निष्काम सेवा आ अमर पद}

\textbf{मूल भोजपुरी रचना:}

\begin{quote}
\addfontfeatures{Color=3D3428}
फल के आस करे जो प्राणी, सो फल पावै अंत बिलाणी।\\
निष्काम भाव से जो रत होई, अमर पदवी पावै सोई॥
\end{quote}

\textbf{भोजपुरी भावार्थ:} जे मनुष्य फल के आशा में कर्म करेला, ओकर फल अंत में खतम हो जाला। बाकिर जे निष्काम भाव से प्रभु भक्ति में लागल रहेला, उहे अमर पद पावेला।

\begin{sloppypar}
{\engfont \textit{English Translation:} One who acts merely for personal rewards reaps results that eventually fade away.
Whoever dedicates life to selfless love and devotion attains immortal peace.}
\end{sloppypar}
